\Article{31}
Whenever the number of quantities \(x,y,z\) etc. is somewhat larger, the numerical determination of the values \(A,B,C\) etc. from the equations \(\xi=0,\allowbreak\eta=0.\quad\zeta=0\) etc. by ordinary elimination becomes rather laborious. For this reason, in Theoria Motus Corporum Coelestium, art.\ 182, we indicated a special algorithm and explained it at length in the \textit{Disquisitione de elementis ellipticis Palladis} (Comm. recent. Soc. Gotting Vol.\ I), by which this labour is made as simple as the matter permits. Namely, the function \(\Omega\) is to be brought into a form such as:
\[
\frac{u^0u^0}{\mathfrak{A}^0}+\frac{u'u'}{\mathfrak{B}'}+\frac{u''u''}{\mathfrak{C}''}+\frac{u'''u'''}{\mathfrak{D}'''}+\text{etc.}+M
\]
Here the divisors \(\mathfrak{A}^0,\mathfrak{B}',\mathfrak{C}'',\mathfrak{D}'''\) etc. are determinate quantities, while \(u^0,u',u'',u'''\) etc. are linear functions of \(x,y,z\) etc.; the second, \(u'\), however, is free of \(x\), the third, \(u''\), free of \(x\) and \(y\), the fourth, \(u'''\), free of \(x,y\) and \(z\), and so on, so that the last, \(u^{(\pi-1)}\), involves only the last of the indeterminates \(x,y,z\) etc. Finally, the coefficients by which \(x,y,z\) etc. are respectively multiplied in \(u^0,u',u''\) etc. are respectively equal to \(\mathfrak{A}^0,\mathfrak{B}',\mathfrak{C}''\) etc. Once this has been done, we must put \(u^0=0,\allowbreak u'=0,\allowbreak u''=0,\allowbreak u'''=0\) etc., from which the values of the unknowns \(x,y,z\) etc. will be obtained most conveniently in reverse order. It seems unnecessary to repeat here the algorithm itself by which this transformation of the function \(\Omega\) is carried out.

But the indeterminate elimination, by means of which the weights of these determinations must be found, requires a much more lengthy calculation still. Inde\SourceLeaf{043}%
ed, the weight of the determination of the last unknown (which alone enters the last \(u^{(\pi-1)}\)) is easily found, from what has been proved in Theoria Motus Corporum Coelestium, to be equal to the last term in the series of divisors \(\mathfrak{A}^0,\mathfrak{B}',\mathfrak{C}''\) etc. Hence several calculators, to avoid that laborious elimination when other means were lacking, resorted to repeating the algorithm of which we have spoken several times, changing the order of the quantities \(x,y,z\) etc. so that each in turn occupied the last place. We therefore hope that geometers will welcome our presenting here a new method of calculating the weights of the determinations, drawn from a deeper investigation of the subject, which seems to leave nothing further to be desired.

\Article{32}
Let us therefore put (I)
\[
\begin{array}{rclllll}
u^0&=&\mathfrak{A}^0x&+\mathfrak{B}^0y&+\mathfrak{C}^0z&+\text{etc.}&+\mathfrak{L}^0\\
u'&=&&\mathfrak{B}'y&+\mathfrak{C}'z&+\text{etc.}&+\mathfrak{L}'\\
u''&=&&&\mathfrak{C}''z&+\text{etc.}&+\mathfrak{L}''\\
\text{etc.}
\end{array}
\]
Then we shall have identically
\[
\begin{aligned}
\tfrac12d\Omega&=\xi\,dx+\eta\,dy+\zeta\,dz+\text{etc.}\\[3pt]
&=\frac{u^0du^0}{\mathfrak{A}^0}+\frac{u'du'}{\mathfrak{B}'}+\frac{u''du''}{\mathfrak{C}''}+\text{etc.}\\[3pt]
&=u^0\left(dx+\frac{\mathfrak{B}^0}{\mathfrak{A}^0}dy+\frac{\mathfrak{C}^0}{\mathfrak{A}^0}dz+\text{etc.}\right)\\
&\qquad+u'\left(dy+\frac{\mathfrak{C}'}{\mathfrak{B}'}dz+\text{etc.}\right)+u''(dz+\text{etc.})+\text{etc.}
\end{aligned}
\]
from which we conclude (II)
\[
\begin{aligned}
\xi&=u^0\\
\eta&=\frac{\mathfrak{B}^0}{\mathfrak{A}^0}u^0+u'\\
\zeta&=\frac{\mathfrak{C}^0}{\mathfrak{A}^0}u^0+\frac{\mathfrak{C}'}{\mathfrak{B}'}u'+u''\\
&\text{etc.}
\end{aligned}
\]
Suppose that from this the following formulas are derived (III)
\SourceLeaf{044}
\[
\begin{aligned}
u^0&=\xi\\
u'&=A'\xi+\eta\\
u''&=A''\xi+B''\eta+\zeta\\
&\text{etc.}
\end{aligned}
\]
Now, from the total differential of the equation
\[
\Omega=\xi(x-A)+\eta(y-B)+\zeta(z-C)+\text{etc.}+M
\]
subtracting the equation
\[
\tfrac12d\Omega=\xi\,dx+\eta\,dy+\zeta\,dz+\text{etc.}
\]
gives
\[
\tfrac12d\Omega=(x-A)d\xi+(y-B)d\eta+(z-C)d\zeta+\text{etc.}
\]
This expression must be identical with the following one derived from III:
\[
\begin{aligned}
&\frac{u^0}{\mathfrak{A}^0}.d\xi+\frac{u'}{\mathfrak{B}'}(A'd\xi+d\eta)\\
&\qquad+\frac{u''}{\mathfrak{C}''}(A''d\xi+B''d\eta+d\zeta)+\text{etc.}
\end{aligned}
\]
From this we conclude (IV)
\[
\begin{array}{rcllll}
x&=&\dfrac{u^0}{\mathfrak{A}^0}&+A'.\dfrac{u'}{\mathfrak{B}'}&+A''.\dfrac{u''}{\mathfrak{C}''}&+\text{etc.}+A\\[7pt]
y&=&&\dfrac{u'}{\mathfrak{B}'}&+B''.\dfrac{u''}{\mathfrak{C}''}&+\text{etc.}+B\\[7pt]
z&=&&&\dfrac{u''}{\mathfrak{C}''}&+\text{etc.}+C\\[4pt]
\text{etc.}
\end{array}
\]
Substituting in these expressions for \(u^0,u',u''\) etc. their values taken from III will complete the indeterminate elimination. In particular, to determine the weights we shall have (V)
\[
\begin{aligned}
\relax[\alpha\alpha]&=\frac1{\mathfrak{A}^0}+\frac{A'A'}{\mathfrak{B}'}+\frac{A''A''}{\mathfrak{C}''}+\frac{A'''A'''}{\mathfrak{D}'''}+\text{etc.}\\[4pt]
[\beta\beta]&=\phantom{\frac1{\mathfrak{A}^0}+{}}\frac1{\mathfrak{B}'}+\frac{B''B''}{\mathfrak{C}''}+\frac{B'''B'''}{\mathfrak{D}'''}+\text{etc.}\\[4pt]
[\gamma\gamma]&=\phantom{\frac1{\mathfrak{A}^0}+\frac1{\mathfrak{B}'}+{}}\frac1{\mathfrak{C}''}+\frac{C'''C'''}{\mathfrak{D}'''}+\text{etc.}\\
&\text{etc.}
\end{aligned}
\]
The simplicity of these formulas leaves nothing to be desired. Equally simple formulas also result for the remaining
\SourceLeaf{045}%
coefficients \([\alpha\beta],[\alpha\gamma],[\beta\gamma]\) etc., but since their use is less frequent, we refrain from giving them here.

\clearpage
